% Part: lambda-calculus
% Chapter: representability
% Section: currying

\documentclass[../../../include/open-logic-section]{subfiles}

\begin{document}

\olfileid[fr]{lam}{rep}{cur}
\olsection{Curryfication}

Une abstraction $\lambd[x][M]$ représente
une fonction d'un seul argument. C'est une
limite lorsque nous voulons définir des
fonctions à plusieurs arguments. Nous
pourrions étendre le lambda-calcul aux
couples, aux triplets, etc. Une fonction
de trois arguments, écrite sous la forme
$\lambd[x][M]$, attendrait alors un triplet
comme argument. La \emph{curryfication} offre
une manière plus commode de procéder.

Considérons un exemple. Supposons un instant
que le lambda-calcul dispose d'une opération
$+$. L'addition prend deux arguments, mais
une abstraction lambda ne définit directement
qu'une fonction d'un seul argument : la
syntaxe n'autorise pas une expression telle
que $\lambd[(x,y)][(x+y)]$. Nous pouvons
toutefois considérer la fonction unaire
$f_x(y)$ donnée par $\lambd[y][(x+y)]$,
qui ajoute~$x$ à son argument~$y$. Il ne
s'agit pas d'une fonction unique, mais
d'une famille de fonctions « ajouter~$x$ »,
une pour chaque nombre~$x$. Définissons
maintenant une autre fonction unaire~$g$
par $\lambd[x][f_x]$. Appliquée à~$x$,
elle renvoie la fonction~$f_x$ : ses
valeurs sont donc elles-mêmes des fonctions.
Si nous appliquons ensuite le résultat à~$y$,
nous obtenons $(g(x))y=f_x(y)=x+y$. Ainsi,
la fonction unaire~$g$ remplit le rôle
de l'addition binaire. La « curryfication »
consiste précisément à représenter une
fonction de deux arguments par une fonction
d'un seul argument dont les valeurs sont
elles-mêmes des fonctions unaires.

Voici maintenant un exemple dans la syntaxe
du lambda-calcul. Comment représenter
$f(x,y)=x$? Une fonction qui reçoit deux
arguments et renvoie le premier s'écrit
$\lambd[x][\lambd[y][x]]$. Elle reçoit
d'abord~$x$ et renvoie la fonction
$\lambd[y][x]$; celle-ci reçoit~$y$,
l'ignore et renvoie toujours~$x$.
Supposons, pour le calcul suivant, que~$y$
ne figure pas libre dans~$M$; sinon,
renommons d'abord cette variable liée.
L'application à deux arguments donne
\begin{align*}
(\lambd[x][\lambd[y][x]])MN
&\bredone (\lambd[y][M])N \\
&\bredone M.
\end{align*}

En général, pour représenter une fonction
d'arguments $x_1$, \dots,~$x_n$ définie
par un terme~$N$, nous pouvons écrire
$\lambd[x_1][\lambd[x_2][\ldots
\lambd[x_n][N]]]$. Soient $M_1$,
\dots,~$M_n$ les arguments. Renommons au
besoin les variables liées $x_1$, \dots,
$x_n$ pour qu'elles soient distinctes et
qu'aucune ne figure libre dans ces arguments.
Posons $N_0\ident N$ et, pour $1\le i\le n$,
\[
N_i\ident\Subst{N_{i-1}}{M_i}{x_i}.
\]
Pour $n\ge2$, la première contraction est
\begin{multline*}
(\lambd[x_1][\lambd[x_2][\ldots
\lambd[x_n][N]]])M_1\dots M_n\bredone{}\\
(\lambd[x_2][\ldots\lambd[x_n][N_1]])
M_2\dots M_n.
\end{multline*}
Après toutes les applications, on obtient
\[
(\lambd[x_1][\lambd[x_2][\ldots
\lambd[x_n][N]]])M_1\dots M_n\red N_n.
\]
Pour $n=1$, ce dernier résultat est
simplement la première contraction. Les
substitutions successives insèrent donc
chaque~$M_i$ à la place de~$x_i$ dans
le corps de la définition, comme le demande
l'application d'une fonction à plusieurs
arguments.

Par exemple, pour appliquer la première projection
à la variable libre $y$ et à un terme $N$,
choisissons un nom $z$ absent des termes en jeu.
Après renommage du second lieur, le calcul est
\[
(\lambd[x][\lambd[z][x]])yN
\bredone (\lambd[z][y])N
\bredone y.
\]
La variable $y$ reste libre tout au long du calcul.

\end{document}
